% Documentación de los datos -- Proyecto integrador. Compilar:
% tectonic -Z search-path=$HOME/repos/ud/maestria/mcic/.UniversidadDistrital-manual-marca/plantillas-latex documentacion_datos.tex
\documentclass[11pt,letterpaper]{article}
\usepackage{ud-informe}
\usepackage{tabularx}
\usepackage{colortbl}
\newcolumntype{L}{>{\raggedright\arraybackslash}X}

% Colores del logo MCIC: cian de acento y azul profundo para títulos
\definecolor{mcicCian}{HTML}{00ADEF}
\definecolor{mcicAzul}{HTML}{005A86}
\colorlet{udrojo}{mcicCian}
\colorlet{udrojoosc}{mcicAzul}
\titleformat{\section}{\cambria\Large\bfseries\color{mcicAzul}}{\color{mcicCian}\thesection.}{0.5em}{}
\titleformat{\subsection}{\cambria\large\bfseries\color{mcicAzul}}{\color{mcicCian}\thesubsection.}{0.5em}{}
\newcommand{\encabezado}{\rowcolor{mcicCian!15}}

\renewcommand{\UDtitulo}{Conociendo los datos}
\renewcommand{\UDsubtitulo}{De dónde salen los datos, cómo están armados y qué tan confiables son}
\renewcommand{\UDtipo}{Ficha técnica del dataset -- Proyecto integrador}
\renewcommand{\UDasignatura}{Inteligencia de Negocios}
\renewcommand{\UDdocente}{Nancy Yaneth Gélvez García}
\renewcommand{\UDestudiantes}{Juan Felipe Rodríguez Galindo -- 20261595004 \\ & Carlos Stiven Mora Hoyos -- 20261595003}
\renewcommand{\UDencabezado}{Proyecto integrador -- Conociendo los datos}

\begin{document}
\UDportada

\section{Origen de los datos}

Los datos vienen del material del caso entregado para el proyecto: un archivo CSV con el histórico
transaccional de una empresa que vende en línea en España y Latinoamérica. Es un dataset \textbf{sintético},
construido para reproducir el comportamiento de un embudo comercial real, y viene acompañado de una evaluación
del caso y de unos parámetros de referencia.

Cada fila es una \textbf{oportunidad comercial}: una persona que llegó por algún canal y que pudo terminar o no
en una compra. El archivo original se conservó sin modificaciones; todo el trabajo se hizo sobre copias y vistas.

\begin{table}[H]
\centering
\caption{Ficha general del dataset.}
\begin{tabular}{ll}
\toprule
\encabezado Característica & Valor \\
\midrule
Formato & CSV, separado por comas \\
Registros & 20.000 oportunidades \\
Columnas & 28 \\
Periodo & 1 de enero de 2024 a 30 de abril de 2026 (28 meses) \\
Ventas & 2.096 (conversión del 10,5\,\%) \\
Ingreso total & €2,33 millones \\
Contribución total & €1,54 millones \\
\bottomrule
\end{tabular}
\end{table}

\section{Estructura}

Las 28 columnas se agrupan en siete bloques:

\begin{table}[H]
\centering\small
\caption{Grupos de variables.}
\begin{tabularx}{\linewidth}{l L}
\toprule
\encabezado Grupo & Variables \\
\midrule
Identificación y tiempo & Identificador, fecha, mes y trimestre \\
Perfil del cliente & Segmento, etapa del ciclo de vida, país y dispositivo \\
Origen & Canal, objetivo de la campaña y nivel de gasto en publicidad \\
Palancas comerciales & Variante de página y de botón, material descargable, pago simplificado, invitación y
asistencia a webinar, oferta del producto nuevo \\
Calificación & Puntaje previo del interesado (1 a 99) \\
Resultado & Compra (sí/no) y días hasta el cierre \\
Valor & Ticket, ingreso, margen, utilidad bruta, costo atribuido y contribución \\
\bottomrule
\end{tabularx}
\end{table}

\begin{table}[H]
\centering\small
\caption{Diccionario de variables.}
\begin{tabularx}{\linewidth}{>{\ttfamily}l l L}
\toprule
\encabezado \normalfont Variable & Tipo & Descripción \\
\midrule
transaction\_id & Identificador & Código único de la oportunidad \\
date, month, quarter & Fecha & Día, mes y trimestre \\
channel & Categórica & Canal de llegada (7 valores) \\
campaign\_objective & Categórica & Objetivo de la campaña (5 valores) \\
customer\_segment & Categórica & Tipo de cliente (5 valores) \\
lifecycle\_stage & Categórica & Visitante nuevo, contacto conocido o cliente recurrente \\
device, geo\_region & Categórica & Dispositivo y país \\
ad\_budget\_level & Ordinal & Nivel de gasto en publicidad (bajo a saturado) \\
campaign\_daily\_spend\_eur & Numérica & Gasto diario de la campaña \\
cost\_attributed\_eur & Numérica & Costo de publicidad atribuido a la oportunidad \\
landing\_variant, cta\_variant & Categórica & Versión de la página y del botón \\
lead\_magnet, checkout\_simplified & Binaria & Material descargable y pago simplificado \\
webinar\_invited, webinar\_attended & Binaria & Invitación y asistencia al webinar \\
new\_product\_offer & Binaria & Se le ofreció el producto nuevo \\
lead\_score & Numérica & Puntaje del interesado (1 a 99) \\
converted\_to\_sale & Binaria & Compró (variable objetivo) \\
days\_to\_close & Numérica & Días hasta la compra \\
aov\_eur, revenue\_eur & Numérica & Ticket e ingreso \\
gross\_margin\_pct, gross\_profit\_eur & Numérica & Margen y utilidad bruta \\
contribution\_profit\_eur & Numérica & Contribución (utilidad después de costos) \\
\bottomrule
\end{tabularx}
\end{table}

\section{Perfilamiento}

\subsection{Distribución de las variables principales}

\begin{table}[H]
\centering\small
\caption{Distribución de las categorías principales.}
\begin{tabularx}{\linewidth}{l L}
\toprule
\encabezado Variable & Distribución (registros) \\
\midrule
Canal & Facebook Ads 4.599 · Google Ads 3.269 · Email 3.254 · SEO/Blog 3.083 · YouTube orgánico 2.391 ·
Afiliados 1.745 · Webinar 1.659 \\
Segmento & B2B Services 4.899 · Ecommerce 4.724 · SMB 4.377 · Creator 4.008 · Enterprise 1.992 \\
País & España 9.228 · México 3.564 · Colombia 1.996 · Argentina 1.978 · Chile 1.652 · EE.\,UU. hispano 1.582 \\
Dispositivo & Escritorio 11.590 · Móvil 7.116 · Tableta 1.294 \\
Ciclo de vida & Visitante nuevo 10.711 · Contacto conocido 6.846 · Cliente recurrente 2.443 \\
Gasto en publicidad & Orgánico 12.132 · Medio 3.290 · Alto 2.029 · Saturado 1.762 · Bajo 787 \\
\bottomrule
\end{tabularx}
\end{table}

El puntaje del interesado va de 1 a 99 con promedio de 56,9. Solo el 10,5\,\% de las oportunidades termina en
compra, así que la variable objetivo está desbalanceada.

\subsection{Conversión por canal}

\begin{table}[H]
\centering\small
\caption{Resumen por canal.}
\begin{tabular}{lrrrr}
\toprule
\encabezado Canal & Registros & Conversión & Ingreso (€) & Contribución (€) \\
\midrule
Email & 3.254 & 19,1\,\% & 680.292 & 505.200 \\
Webinar & 1.659 & 17,5\,\% & 355.556 & 253.968 \\
SEO / Blog & 3.083 & 12,2\,\% & 420.311 & 300.375 \\
YouTube orgánico & 2.391 & 10,8\,\% & 284.802 & 207.307 \\
Afiliados & 1.745 & 7,3\,\% & 136.072 & 78.113 \\
Google Ads & 3.269 & 6,1\,\% & 221.094 & 111.850 \\
Facebook Ads & 4.599 & 4,8\,\% & 228.624 & 84.433 \\
\bottomrule
\end{tabular}
\end{table}

Los canales pagados traen más volumen pero convierten menos; los propios (email, webinar, contenido) convierten
más y dejan más contribución.

\section{Calidad de los datos}

\subsection{Validaciones}

Se revisaron cuatro frentes y \textbf{ninguna regla presentó fallas}:

\begin{itemize}
  \item \textbf{Unicidad y completitud}: no hay identificadores repetidos ni nulos en los campos clave.
  \item \textbf{Coherencia temporal}: el mes y el trimestre coinciden con la fecha.
  \item \textbf{Dominios}: las binarias solo toman 0 y 1, el puntaje está entre 1 y 99 y el margen entre 0 y 1.
  \item \textbf{Coherencia entre columnas}: nadie asiste a un webinar sin invitación, los días de cierre solo
        existen si hubo venta, el ingreso cuadra con el ticket y el margen, y solo los canales pagados tienen
        nivel de gasto en publicidad.
\end{itemize}

\subsection{Hallazgos}

Aunque el dato está limpio, el perfilamiento mostró particularidades que había que tener en cuenta:

\begin{table}[H]
\centering\small
\caption{Hallazgos del perfilamiento y tratamiento.}
\begin{tabularx}{\linewidth}{>{\raggedright\arraybackslash}p{3.3cm} L L}
\toprule
\encabezado Hallazgo & Descripción & Tratamiento \\
\midrule
Costo fijo oculto & En las 2.096 ventas la contribución es €22 menor de lo esperado & Se documentó como
costo fijo por venta \\
Gasto diario no sumable & El gasto diario de campaña cambia en cada fila; sumarlo da una cifra sin sentido &
Solo se promedia; para sumar se usa el costo atribuido \\
Ceros estructurales & Sin venta, el ticket vale 0 y los días de cierre quedan vacíos (17.904 filas) & El
ticket se calcula solo sobre ventas y los vacíos no se imputan \\
Variables posteriores a la venta & Ingreso, margen, contribución y días de cierre solo se conocen después &
No se usan para predecir la compra \\
Producto nuevo limitado & Solo se ofreció en tres segmentos (B2B Services, Ecommerce y Enterprise) & Se compara
dentro de cada segmento \\
Último trimestre incompleto & 2026T2 solo tiene abril & Se marca como parcial en las comparaciones \\
\bottomrule
\end{tabularx}
\end{table}

\section{Conclusión}

El dataset es completo y consistente: 20.000 oportunidades en 28 meses, sin nulos ni duplicados en los campos
clave y sin fallas en las reglas de validación. Las particularidades encontradas (costo fijo, gasto no sumable,
ceros estructurales y variables posteriores a la venta) no son errores, sino características que se deben
respetar al calcular indicadores y al construir modelos.

\end{document}
